% TOP_PROBLEM: {"id": 1101002, "title": "Questions in Geometric Group Theory — Q 10.2", "category_id": 17, "category": {"id": 17, "name": "group_theory", "display_name": "Group Theory", "description": "Problems about groups, group actions, representations, and related algebraic structures.", "slug": "group-theory", "order_index": 17, "created_at": "2026-07-31T15:26:25.671Z"}, "status": "open", "problem_number": "AMR-010-1002", "set_id": 13, "statusQualification": "Retains the source general CW-complex domain, without adding finiteness or asphericity."}
% FIRST_POSTED: 2026-10-01
% ENTRY_KIND: statement_only
% UnsolvedMath ID 1101002; source catalogue code AMR-010-1002
% !TeX program = lualatex
% Definition and sources only; this file is not an LLM solution attempt.
\documentclass[11pt,a4paper]{article}
\usepackage[margin=25mm]{geometry}
\usepackage{fontspec}
\setmainfont{lmroman10-regular.otf}[BoldFont=lmroman10-bold.otf,
 ItalicFont=lmroman10-italic.otf,BoldItalicFont=lmroman10-bolditalic.otf]
\usepackage{amsmath,amssymb,mathtools}
\usepackage{unicode-math}
\setmathfont{latinmodern-math.otf}
\AtBeginDocument{\let\setminus\smallsetminus}
\usepackage{xurl,hyperref}
\hypersetup{colorlinks=true,urlcolor=blue}
\setcounter{secnumdepth}{0}
\setlength{\parindent}{0pt}
\setlength{\parskip}{5pt}
\setlength{\emergencystretch}{3em}
\begin{document}
\section*{Questions in Geometric Group Theory — Q 10.2}\label{1101002}
{\small UnsolvedMath ID 1101002 · AMR-010-1002 · Group Theory · AMR Open Problem Lists}
\subsection{Definitions and mathematical statement}
Let $X$ be a connected CW complex with no cells of dimension greater than three, and put $\pi=\pi_1(X)$. A local coefficient system of abelian groups on $X$ consists of groups transported isomorphically along paths, with transport depending only on endpoint-preserving homotopy and respecting composition. After choosing a basepoint, such a system is equivalent to a $\mathbb Z\pi$-module, with monodromy given by the fundamental-group action.

Its cohomology $H^3(X;\mathcal M)$ can be computed from the cellular chains of the universal cover using the corresponding module-valued cochains. Assume
\[
 H^3(X;\mathcal M)=0
 \qquad\text{for every local coefficient system }\mathcal M.
\]
Does there exist an $X$ satisfying this condition which is not homotopy equivalent to any two-dimensional CW complex?

The target two-complex is allowed to have infinitely many cells, and the source wording does not require $X$ to be finite. The finite-complex version is an important additional restriction, but it is not silently substituted here. No asphericity assumption is made; higher homotopy in degree two can be present. Vanishing only for constant integer coefficients would be much weaker than the universal condition displayed above. Likewise, a particular three-cell presentation that cannot be collapsed directly does not constitute a counterexample, because an unrelated two-complex could still have the same homotopy type. The problem concerns homotopy dimension rather than a chosen cell decomposition.
\subsection{Short English statement}
Can a three-dimensional complex have vanishing third cohomology for every local system yet fail to have the homotopy type of a two-complex?
\subsection{Sources}
\begin{itemize}
\item[S1] ULAM.AI and the UnsolvedMath Contributors, supplied problems.json, record 1101002 (AMR-010-1002), AMR Open Problem Lists. Catalogue identity and source transcription; CC BY 4.0.
\url{https://huggingface.co/datasets/ulamai/UnsolvedMath}.
\item[S2] Bestvina - Questions in Geometric Group Theory (pdf) (2004), Question 10.2 (PDF page 18). Original source indexed by AMR-010-1002.
\url{https://www.math.utah.edu/~bestvina/eprints/questions-updated.pdf}.
\end{itemize}
\end{document}
